\documentclass[10pt,a4paper]{amsart}
\usepackage[margin=19mm]{geometry}
\usepackage{amsmath,amssymb,mathtools}
\usepackage[scr=rsfs,cal=euler]{mathalfa}
\usepackage{xfrac}
\usepackage[unicode,hidelinks,bookmarks=false]{hyperref}
\newcommand{\ie}{i.e.\ }
\newcommand{\cf}{cf.\ }
\newcommand{\resp}{respectively\ }
\newcommand{\Subsection}{\subsection}
\providecommand{\nobreakdash}{\nobreak\hbox{-}}
\setlength{\emergencystretch}{2em}
\multlinegap0pt
\usepackage{amssymb}
\usepackage{tikz}
\newtheorem{lem}{Lemma}
\def\Ric {\mathrm{Ric}}
\def\d {\mathrm{d}}
\newcommand{\dv}{\operatorname{div}}
\def\la{\lambda}
\renewcommand{\leq}{\leqslant}
\def\na{\ensuremath{\nabla}}
\newcommand{\norm}[1]{\left\lVert#1\right\rVert}%norm%
\title{Liouville theorem on $p$-biharmonic map from gradient Ricci soliton}
\author{Xiangzhi Cao}
\date{Source: arXiv:2603.16541v2 (CC BY 4.0)}
\begin{document}
\maketitle

\noindent\emph{Source and licence.} Liouville theorem on $p$-biharmonic map from gradient Ricci soliton. Version arXiv:2603.16541v2 (CC BY 4.0), distributed by arXiv under the Creative Commons Attribution 4.0 International licence, http://creativecommons.org/licenses/by/4.0/, as recorded in the arXiv licence metadata for this exact version and shown on the abstract page https://arxiv.org/abs/2603.16541v2. Full licence notice: \texttt{licenses/arxiv-2603.16541-CC-BY-4.0.txt}.\par\smallskip

\noindent\textit{Editorial scope.} Counted unit: the article's lem 15 (source0.tex lines 2314--2331). The article's complete original proof of it is delivered with it (source0.tex lines 2332--2359, one \texttt{proof} environment of 28 lines). No mathematics is written, repaired or re-derived: the statement and the proof are the article's own lines, in the article's own notation, and no retained line is substituted or re-flowed.  Retained uncounted as the source-local context the proof needs: lines 2185--2197.  Nothing was omitted from any retained span, so the package carries no omission record.  No retained line contains a citation, so the wrapper carries no bibliography and no bibliography entry.  No retained line contains a figure environment, an image-inclusion command or drawing code, so no image is delivered and the package declares no asset dependency.  Editorial formatting only, in addition: an AMS \texttt{amsart} preamble that restates the article's own declarations and macros used by the retained lines, namely the environments \texttt{lem} on separate counters, together with the standard packages those lines need.  The retained environments print in this document as Lemma 1 in order of appearance, whereas the article prints the same items with the numbering of its own full document.  The complete native source of the article as distributed in the arXiv e-print for this exact version is bundled unchanged as \texttt{sources/196562/source0.tex}.
\begin{equation}\label{13}
	\begin{split}
			&\int_M\eta^2\big(\lambda(m-4)-\operatorname{Scal}^M\big)\|\tau_p(\phi)\|^2\dv\\
			&+4\int_M\eta^2\operatorname{Ric}^M(e_i,e_j)\langle\na \left(\norm{d\phi}^{p-2} d\phi \right)(e_i,e_j),\tau_p(\phi)\rangle \dv \\
		&+\int_M 2(p-2)\norm{\tau_p(\phi)}^{p-4}\left\langle \d \phi,\na \tau_p(\phi) \right\rangle\left\langle \d \phi(e_i),\d \phi(e_j) \right\rangle\Ric_{ij}\\&-\int_M \eta^2(p-2)\la\norm{\tau_p(\phi)}^{p-4}\left\langle \d \phi,\na \tau_p(\phi) \right\rangle\norm{\na \phi}^2 \dv\\
		=&
		-4\int_M\norm{\d \phi}^{p-2}\langle d\phi(\operatorname{Ric}^M(\nabla\eta^2)),\tau_p(\phi)\rangle \dv 
		+4\lambda\int_M\langle d\phi(\nabla\eta^2),\tau_p(\phi)\rangle \dv 
		\\
		&-\int_M\langle\nabla\eta^2,\nabla f\rangle\|\tau_p(\phi)\|^2 \dv +2\int_M(\Delta\eta^2)\norm{\d \phi}^{p-2}\langle d\phi(\nabla f),\tau_p(\phi)\rangle \dv 
		\\&+4\int_M(\nabla_{e_i}\eta^2)\nabla_{e_j}f\norm{\d \phi}^{p-2}\langle\na d\phi(e_i,e_j),\tau_p(\phi)\rangle \dv,
	\end{split}
\end{equation}

\begin{lem}\label{15}
	Assume that $(M,g,f)$ is a complete, non-compact gradient Ricci soliton
	with $\operatorname{Ric}^M\leq C$ and $\|\nabla f\|<\infty$.
	Let $\phi\colon M\to N$ be a smooth $p$-biharmonic map with 
	\begin{equation}\label{14}
		\int_M(\|d\phi\|^{4p-6}+\|\na d\phi\|^2)\dv+\int_M\norm{\d \phi}^{2p-4} \|\na d\phi\|^2 \dv<\infty.
	\end{equation}
	Then the following inequality holds
	\begin{equation}\label{20}
		\begin{split}
			&\int_M\big(\lambda(m-4)-\operatorname{Scal}^M\big)\|\tau_p(\phi)\|^2\dv
			+4\int_M\operatorname{Ric}^M(e_i,e_j)\langle\na d\phi(e_i,e_j),\tau_p(\phi)\rangle \dv\\
			&+\int_M (p-2)\norm{\tau_p(\phi)}^{p-4}\left\langle \d \phi,\na \tau_p(\phi) \right\rangle\left\langle \d \phi(e_i),\d \phi(e_j) \right\rangle\Ric_{ij}\\
			&+\int_M (p-2)\norm{\tau_p(\phi)}^{p-4}\left\langle \d \phi,\na \tau_p(\phi) \right\rangle\left\langle \d \phi(e_i),\d \phi(e_j) \right\rangle\Ric_{ij}\\&+\int_M \eta^2(p-2)\norm{\tau_p(\phi)}^{p-4}\left\langle \d \phi,\na \tau_p(\phi) \right\rangle\norm{\na \phi}^2\bigg) \dv\\
			&\leq 0.
		\end{split}
	\end{equation}
\end{lem}

\begin{proof}[Proof of Lemma \ref{15}]
	Again, let  $0\leq\eta\leq 1$ on $M$ be such that
	\begin{equation*}
		\eta(x)=1\textrm{ for } x\in B_R(x_0),\qquad \eta(x)=0\textrm{ for } x\in B_{2R}(x_0),\qquad |\nabla^q\eta|\leq\frac{C}{R^q}\textrm{ for } x\in M,
	\end{equation*}
	where $B_R(x_0)$ denotes the geodesic ball around the point $x_0$ with radius $R$ and $q=1,2$.
	Moreover, let $\{e_i\},i=1,\ldots,m$ be an orthonormal basis of $TM$ that satisfies 
	$\nabla_{e_i}e_j=0,i,j=1,\ldots,m$ at a fixed point $p\in M$.
	
	We need to estimate the terms on 	the right hand side of \eqref{13}. 	By the  assumption \eqref{14},	it is easy to infer
	\begin{align*}
		&\int_M \norm{\d \phi}^{p-2}\langle d\phi(\operatorname{Ric}^M(\nabla\eta^2)),\tau_p(\phi)\rangle \dv\leq\frac{C}{R}\int_M(\|\na d\phi\|^2+\|d\phi\|^{4p-6})\dv \to 0, \\
		&\int_M \norm{\d \phi}^{p-2} \langle d\phi(\nabla\eta^2),\tau_p(\phi)\rangle \dv 
		\leq\frac{C}{R}\int_M(\|\na d\phi\|^2+\|d\phi\|^{4p-6})\dv 
		\to 0,  \\
		&\int_M\langle\nabla\eta^2,\nabla f\rangle\|\tau_p(\phi)\|^2 \dv 
		\leq\frac{C}{R}\int_M \norm{\d \phi}^{2p-4}\|\na d\phi\|^2\dv 
		\to 0, \\
	&\int_M \norm{\d \phi}^{p-2}(\Delta\eta^2)\langle d\phi(\nabla f),\tau_p(\phi)\rangle \dv 
		\leq\frac{C}{R^2}\int_M(\|\na d\phi\|^2+\|d\phi\|^2)\dv 
		\to 0, \\
		&\int_M\norm{\d \phi}^{p-2}(\nabla_{e_i}\eta^2)\nabla_{e_j}f\langle\na d\phi(e_i,e_j),\tau_p(\phi)\rangle 
		\leq\frac{C}{R}\int_M\norm{\d \phi}^{2p-4} \|\na d\phi\|^2 \dv\to 0
	\end{align*}
	as $R\to\infty$.
	Note that we used the inequality $\|\tau_p(\phi)\|\leq C_p \norm{\d \phi}^{p-2}\|\na d\phi\|$ and $ \Ric\leq C, \norm{\na f}\leq C$.
	This completes the proof.
\end{proof}

\end{document}
